La hipótesis tiene nombre: compartición falsa.

Recordad que la caché no mueve datos sueltos, sino líneas de 64 bytes. En bench.cpp las pilas están seguidas en memoria. Cada pila de Autores ocupa 8 bytes, así que caben ocho en una línea. Las diez pilas ocupan unas dos líneas.

El dibujo lo muestra: A, B y C escriben en pilas distintas, pero en la misma línea. Para escribir, cada uno necesita la línea entera en exclusiva, así que viaja de uno a otro como si compartieran un dato. Es como diez personas con su propio cuaderno, pero los diez cuadernos son renglones de la misma hoja.

El programa oficial pide 128 bytes por pila, con alignas: cada pila en su hueco. Curiosamente, el diseño compacto de Autores, 8 bytes frente a 16 de los demás, aquí le perjudica.
